\subsection{Висновки до першого розділу}

У цьому розділі було проведено комплексний огляд літератури та глибокий теоретичний аналіз математичних основ просторового масштабування цифрових зображень. На базі дослідженого матеріалу було фундаментально обґрунтовано вибір оптимального інтерполяційного апарату для подальшої алгоритмічної реалізації. Головні результати розділу можна звести до наступних ключових положень:

\begin{itemize}
	\item \textbf{Сформовано детерміновану математичну модель цифрового зображення.} Перехід від фізичного неперервного поля до нормалізованого двовимірного масиву пікселів (та його векторно-матричного подання) дозволив звести задачу масштабування до строгих операцій лінійної алгебри. Також класифіковано основні візуальні артефакти масштабування, такі як феномен Гіббса (ефект дзвону, перерегулювання), аліасинг та розмиття, які є прямим наслідком математичних обмежень методів інтерполяції.
	
	\item \textbf{Доведено неспроможність використання рівномірних сіток для глобальної інтерполяції.} Аналітично та візуально підтверджено, що застосування поліномів високого ступеня на класичній рівновіддаленій матриці пікселів викликає феномен Рунге. Це призводить до експоненціального накопичення похибки та руйнування візуальної інформації на краях зображення.
	
	\item \textbf{Обґрунтовано перехід до нерівномірної топології дискретизації.} Математично доведено, що використання екстремумів поліномів Чебишева (вузлів Чебишева) як точок просторової сітки дозволяє повністю нівелювати феномен Рунге завдяки асимптотичному згущенню вузлів на краях інтервалу, забезпечуючи рівномірну збіжність інтерполянта.
	
	\item \textbf{Виведено та оптимізовано математичний апарат барицентричної інтерполяції.} Доведено, що застосування другої барицентричної формули є найбільш ефективним алгоритмічним рішенням. Для вузлів Чебишева вагові коефіцієнти цієї формули зводяться до елементарного знакозмінного ряду, що знижує обчислювальну складність, гарантує масштабну інваріантність і забезпечує абсолютну числову (пряму) стійкість навіть для зображень із високою роздільною здатністю.
	
	\item \textbf{Доведено високу швидкість збіжності для реальних зображень.} На основі апроксимаційної теореми Вейєрштрасса та аналізу коефіцієнтів Чебишева встановлено, що навіть для негладких функцій (якими є цифрові зображення з різкими контурами та межами об'єктів), розроблений метод гарантує стабільну алгебраїчну збіжність $O(n^{-\nu})$. Це дозволяє локалізувати похибку і зберігати структурну цілісність контурів без експоненціального розкачування.
\end{itemize}

Таким чином, викладений теоретичний базис переконливо доводить, що комбінація вузлів Чебишева та другої барицентричної формули є математично надійною основою для розробки високоякісного алгоритму масштабування, практична реалізація та оптимізація якого розглядається у наступному розділі.